\clearpage
\VerdierPhysicalPageAnchor{0074}%
\VerdierSemanticAnchor{ch1-2-4-1-pi-opposite}
它是 (2.2.3) 和 (2.2.4) 中所定义函子的反函子。同样，记：

\VerdierSemanticAnchor{ch1-2-4-1-eq4}
\begin{equation}
  \int_{\varphi,C^\circ}^{\Pi}\colon
  \operatorname{comp}^{n}(C^\circ)
  \longrightarrow
  \operatorname{comp}^{m}(C^\circ),
  \tag{2.4.1.4}
\end{equation}

\VerdierSemanticAnchor{ch1-2-4-1-pi-definition}
为 (2.2.3) 和 (2.2.4) 中所定义的函子。直接验证可知，图式：

\VerdierSemanticAnchor{ch1-2-4-1-eq5}
\begin{equation}
  \resizebox{0.82\textwidth}{!}{%
    $\begin{tikzcd}[ampersand replacement=\&,column sep=6.5em,row sep=5.2em]
      \operatorname{comp}^{n}(C)^\circ
        \arrow[r,"\iota_n"]
        \arrow[d,"{\left(\int_{\varphi,C}^{\Sigma}\right)^\circ}"']
        \& \operatorname{comp}^{n}(C^\circ)
          \arrow[d,"{\int_{\varphi,C^\circ}^{\Pi}}"] \\
      \operatorname{comp}^{m}(C)^\circ
        \arrow[r,"\iota_m"]
        \& \operatorname{comp}^{m}(C^\circ)
    \end{tikzcd}$%
  }
  \tag{2.4.1.5}
\end{equation}

\VerdierSemanticAnchor{ch1-2-4-1-eq5-commutative}
是交换的。特别地，当 $\varphi$ 是唯一的映射
$[n]\longrightarrow[1]$ 时，有交换图式：

\VerdierSemanticAnchor{ch1-2-4-1-eq6}
\begin{equation}
  \resizebox{0.82\textwidth}{!}{%
    $\begin{tikzcd}[ampersand replacement=\&,column sep=6.5em,row sep=5.2em]
      \operatorname{comp}^{n}(C)^\circ
        \arrow[r,"\iota_n"]
        \arrow[d,"{\left(\int_{n,C}^{\Sigma}\right)^\circ}"']
        \& \operatorname{comp}^{n}(C^\circ)
          \arrow[d,"{\int_{n,C^\circ}^{\Pi}}"] \\
      \operatorname{comp}(C)^\circ
        \arrow[r,"\iota_1"]
        \& \operatorname{comp}(C^\circ)
    \end{tikzcd}$%
  }
  \tag{2.4.1.6}
\end{equation}

\VerdierSemanticAnchor{ch1-2-4-1-twisted-isomorphism}
并且可立即验证，同构对 $(\iota_1,\iota_n)$ 是扭曲
$\mathbb Z[n]$-函子 $\left(\int_{n,C}^{\Sigma}\right)^\circ$ 到扭曲
$\mathbb Z[n]$-函子 $\int_{n,C^\circ}^{\Pi}$ 的一个同构
（$\int_{n,C^\circ}^{\Pi}$ 的分次结构在 (2.3.4) 中描述；
$\left(\int_{n,C}^{\Sigma}\right)^\circ$ 的分次结构在 (2.3.4) 和
(1.4.8) 中描述）。

\VerdierSemanticAnchor{ch1-2-4-1-identification}
今后，我们通过同构 $\iota_n$ 将范畴
$\operatorname{comp}^{n}(C)^\circ$ 与
$\operatorname{comp}^{n}(C^\circ)$，以及函子
$\left(\int_{\varphi,C}^{\Sigma}\right)^\circ$ 与
$\int_{\varphi,C^\circ}^{\Pi}$ 相互识别。

\VerdierSemanticAnchor{ch1-2-4-2}
\noindent\textbf{2.4.2.}\quad
设 $C_i$（$i\in[n]$）为一族加法范畴，
$\prod_{i\in[n]}C_i$ 为它们的积。以 $(X_1,X_2,\ldots,X_n)$ 表示
$\prod_{i\in[n]}C_i$ 的对象，并使用记号
$(X_1,X_2,\ldots,f_i,\ldots,X_n)$ 表示态射：

\VerdierSemanticAnchor{ch1-2-4-2-component-morphism}
\begin{equation*}
  (\operatorname{id}_{X_1},\operatorname{id}_{X_2},\ldots,
  f_i,\ldots,\operatorname{id}_{X_n}).
\end{equation*}
